%-----------------------------------------------------------------
%Di sini awal masukan untuk Prakata
%------------------------------------------------------------
\preface
\justifying
\noindent Puji syukur penulis panjatkan ke hadirat Allah SWT atas berkah dan rahmat-Nya sehingga skripsi ini dapat terselesaikan dengan baik. Skripsi ini dibuat untuk menyelesaikan pendidikan jenjang sarjana pada Institut Teknologi Sumatera. Penyusunan skripsi ini banyak mendapat bantuan dan dukungan dari berbagai pihak sehingga dalam kesempatan ini, dengan penuh kerendahan hati, penulis mengucapkan terima kasih kepada:
\begin{enumerate}
\item{Prof. Dr. I Nyoman Pugeg Aryantha selaku Rektor Institut Teknologi Sumatera,}
\item{Dr. Ikah Ning Prasetiowati Permanasari, S.Si., M.Si. selaku Dekan Fakultas Sains Institut Teknologi Sumatera,}
\item{Tirta Setiawan, S.Pd., M.Si. selaku Koordinator Program Studi,}
\item{Luluk Muthoharoh, M.Si. selaku dosen wali akademik dan juga pembimbing penelitian yang telah memberikan arahan dan bimbingan selama perkuliahan,}
\item{Ardika Satria, M.Si. selaku dosen pembimbing skripsi yang telah banyak memberikan bimbingan, saran, serta motivasi kepada penulis dalam penyusunan skripsi ini,} 
\item{Teman-teman seperjuangan selama perkuliahan di Institut Teknologi Sumatera, khususnya teman-teman angkatan 2022 Program Studi Sains Data,}
\end{enumerate}

Penulis menyadari bahwa penyusunan skripsi ini jauh dari sempurna.
Akhir kata penulis mohon maaf yang sebesar-besarnya apabila ada kekeliruan di dalam penulisan skripsi ini.



\vspace{0.5cm}

\begin{flushright}
\begin{tabular}{p{3.5cm}l}
&Lampung Selatan, \approvaldatenc \\[1.5cm]
&\textbf{\fullnamenc}
\end{tabular}
\end{flushright}